% ---- pg-555 （书页 543） ----

离(相似性)矩阵, 也不必存储基本数据, 所以同分层方法相比, 非分层方法可用于大得多的数据集.

非分层方法的起点或者是(1)一个初始的分割, 或者是(2)构成未来聚类核心的一组初始种子点. 好的初始结构应该没有明显的偏倚.
一种起动方式是从项目中随机选出种子点, 或者将项目随机地分割成初始组的集合.

在这一节中, 我们将讨论一种比较流行的非分层方法——$K$ 均值法.

\mvsub{$K$ 均值法}

麦奎因(MacQueen, 见文献[25])将他所提出的一种算法命名为 $K$ 均值法. 这种算法的基本思想是将每一个项目分给具有最近中心(均值)的聚类.
最简单形式中的算法包括以下三个步骤:

\noindent 1.\quad 将所有项目分成 $K$ 个初始聚类.

\noindent 2.\quad 将项目表中的某个项目划入中心(均值)离它最近的聚类. (这里的距离通常是用标准化或非标准化数据算出的欧氏距离.) 对得到项目和失去项目的两个聚类重新计算它们的中心坐标.

\noindent 3.\quad 重复步骤 2, 直至所有的点都不能再分配时为止. \hfill (12-16)

\noindent 在上述过程中, 步骤 1 也可以不从分割出 $K$ 个初始聚类开始, 而从规定 $K$ 个初始中心(种子点)开始, 然后进入步骤 2.

最终的分类结果在某种程度上依赖于初始分组或初始种子点的选择. 经验显示, 聚类过程中的绝大多数重要变化均发生在第一次再分配中.

\noindent\textbf{例 12.11}\quad (用 $K$ 均值法聚类)

假定我们对 $A,B,C,D$ 四个项目分别测量两个变量 $x_1$ 和 $x_2$, 得到如下数据:

\begin{mvtable}
\footnotesize
\begin{tabular}{lcc}
\toprule
\multirow{2}{*}{项\quad 目} & \multicolumn{2}{c}{变\quad 量}\\
\cmidrule(lr){2-3}
 & $x_1$ & $x_2$\\
\midrule
$A$ & 5 & 3\\
$B$ & $-1$ & 1\\
$C$ & 1 & $-2$\\
$D$ & $-3$ & $-2$\\
\bottomrule
\end{tabular}
\end{mvtable}

我们的目标是将这些项目分成 $K=2$ 个聚类, 使每个聚类内部的项目之间的距离比与分别属于不同聚类的项目之间的距离小.
为了实施 $K=2$ 均值法, 我们将这些项目先随意分成两个聚类, 比如说 $(AB)$ 和 $(CD)$, 然后计算这两个聚类的中心(均值)的坐标 $(\bar{x}_1,\bar{x}_2)$.
这样, 在步骤 1 中我们有:

\begin{mvtable}
\footnotesize
\begin{tabular}{lcc}
\toprule
\multirow{2}{*}{聚\quad 类} & \multicolumn{2}{c}{中心的坐标}\\
\cmidrule(lr){2-3}
 & $\bar{x}_1$ & $\bar{x}_2$\\
\midrule
$(AB)$ & $\dfrac{5+(-1)}{2}=2$ & $\dfrac{3+1}{2}=2$\\[6pt]
$(CD)$ & $\dfrac{1+(-3)}{2}=-1$ & $\dfrac{-2+(-2)}{2}=-2$\\
\bottomrule
\end{tabular}
\end{mvtable}
